\chapter{SAM tools}\label{chap:samtools}

The SAM tools read, write, aggregate, balance and analyse a SAM before it is
mapped and calibrated. They change the SAM only; the model is unchanged. Each
command that changes the SAM starts a new model from the result, as
\manref{cge_sam} does: earlier mappings, calibrations and simulations are
discarded. In this chapter $\bar X$ is the SAM before a command and $X$ the SAM
after it, with $x_{ij}$ the payment from column account $j$ to row account
$i$.

\section{Reading a SAM}\label{sec:reading}

\manref{cge_sam} reads a SAM in one of three forms:
\begin{enumerate}
\item \textbf{Wide data}: one observation per account, a string variable with
the account labels and one numeric variable per account, from data in memory or
from a file read with \cmd{use}, \cmd{import delimited} or \cmd{import excel}.
\item \textbf{A wide CSV file read directly}: the file is parsed in Mata and
the data in memory are not changed. This is automatic when the SAM has more
accounts than Stata can hold as variables (\cmd{c(maxvar)}; 2,048 in
Stata/IC), and is chosen for any CSV file with the option \cmd{direct}.
\item \textbf{Long data}: one observation per non-zero cell, with string
variables for the row and column accounts and a numeric variable for the
value. The number of accounts is limited only by memory. Accounts are ordered
as they first appear as row accounts, followed by any that appear only as
column accounts, so that a SAM written in long form and read again keeps its
order. A cell given twice is an error.
\end{enumerate}
In the direct reader the first line holds the column labels, the first field of
each following line the row label, and the remaining fields the values. Fields
may be quoted. A blank field is a zero, counted in \res{r(N_blank)}; a field
that is not a number (for example \cmd{1,234} with a thousands separator) is an
error, never a zero. Row and column labels must agree, in the same order.

\textbf{Labels.} Account labels are used as names in \manref{cge_map}, so they
must be valid Stata names. In wide data, Stata renames invalid labels when the
file is imported (\cmd{s-i} becomes \cmd{si}), and the variable names are
used. In the direct reader and in long data the same rule is applied:
characters other than letters, digits and the underscore are removed, then any
leading digits, and the result is cut to 32 characters. If a label cannot be
converted (the result is empty, a reserved word, or the same as another
label), \cmd{cge_sam} stops and names it. Each change is displayed.

\section{Writing a SAM}\label{sec:writing}

\manref{cge_samexport} writes the SAM of the current model, as read, aggregated
or balanced, to a \cmd{.csv}, \cmd{.dta} or \cmd{.xlsx} file, or loads it as
data. Wide output has one variable per account (\cmd{.dta} and \cmd{.xlsx}
need the accounts to fit as variables; \cmd{.csv} is written in Mata and has no
such limit). Long output has the variables \cmd{row}, \cmd{column} and
\cmd{value}, one observation per non-zero cell. With \cmd{mapped}, the SAM after
the layout rewrites of section~\ref{sec:layouts} is written instead, with the
internal tax accounts. Numbers are written in the shortest of \texttt{\%21.0g}
and \texttt{\%24.17g} that reads back as the same double, so that a SAM written
and read again is unchanged.

\section{Mapping templates}\label{sec:template}

\manref{cge_template} writes a template with one row per account and the
columns \cmd{account}; \cmd{role}, the \cmd{cge_map} options that list the
account; \cmd{pair}, for an activity, the commodity it is paired with;
\cmd{group}, for aggregation, the account it is merged into; and the account's
row and column totals and their difference. An Excel template also has a sheet
\cmd{roles} that lists the valid entries. When the model in memory is already
mapped, \cmd{role} and \cmd{pair} are filled from the mapping.

A role is any option name of \cmd{cge_map}, in either case, with the singular
forms (\cmd{activity}, \cmd{commodity}, \cmd{factor}, \cmd{household},
\cmd{enterprise}, \cmd{margin}, \cmd{marginaccount}, \cmd{stock}) also
accepted. An account may have several roles when \cmd{cge_map} allows it, for
example \cmd{activities margins} for a margin activity, or
\cmd{government saltax} when sales taxes are paid directly to the government.
A blank role leaves the account unmapped, which \cmd{cge_map} accepts only for
accounts without flows.

\cmd{cge_map using} reads a template. Each option lists its accounts in the
order of the SAM. Activities are paired with commodities by the \cmd{pair}
column; a blank \cmd{pair} pairs the activity with the commodity whose column
pays it most (the largest cell $x_{ac}$). A tie, no sales to any commodity, or a
commodity paired with two activities is an error. After aggregation with the
same template, an account of the current SAM may be a group: it takes the role
of its members, which must agree, and is paired with the group of its members'
pairs. The options that result are displayed, so that the same mapping can be
written as a command.

\section{Aggregation}\label{sec:aggregation}

\manref{cge_aggregate} merges accounts into groups. With the indicator matrix
$G$ of groups by accounts ($G_{gk} = 1$ if account $k$ is in group $g$),
\begin{equation*}
X = G\,\bar X\,G',
\end{equation*}
so that every cell is the sum of the cells it merges. Totals add up, and an
account's balance is the sum of its members' balances. A group takes the
position of its first member. Flows among the members of a group become a
diagonal cell of the new account; they are kept, as transfers within the
group. The model accepts diagonal cells for households, enterprises and the
government, and \cmd{cge_map} reports any other as a flow the model does not
represent. The option \cmd{netdiagonal} removes them; totals then fall by the
diagonal cell, and balances do not change.

\section{Balancing}\label{sec:balancing}

\manref{cge_balance} makes the SAM balanced, every account's row total equal to
its column total. Both of its methods keep the sign of every cell, every zero
cell, and the cells listed in \cmd{fixed()} at their values in $\bar X$. Both
minimise the same measure of the information added to $\bar X$
\citep{junius2003,lenzen2007}, over the cells that may change,
\begin{equation*}
D(X, \bar X) = \sum_{ij} \Bigl(|x_{ij}|\,\ln\frac{|x_{ij}|}{|\bar x_{ij}|} - |x_{ij}| + |\bar x_{ij}|\Bigr),
\end{equation*}
which for non-negative matrices is the cross entropy that RAS minimises
\citep{mcdougall1999}. They differ in what is known about the totals.

\subsection{RAS}\label{sec:ras}

Every account $k$ has a target total $t_k$, and the row and column totals of
$X$ are set to it. By default $t_k$ is the average of the account's row and
column totals in $\bar X$; the targets may instead be the row or the column
totals, or be given for individual accounts. If an account's row has no cell
that may change, its target is its row total (and likewise for its column); if
both are fixed and differ, or a given target contradicts them, no balanced SAM
exists and \cmd{cge_balance} stops.

With $P$ the positive part and $N$ the absolute value of the negative part of
the cells that may change, and $u$ and $v$ the targets less the fixed cells of
each row and column, generalised RAS (GRAS) finds row factors $r$ and column
factors $s$ such that
\begin{equation*}
x_{ij} = r_i\,p_{ij}\,s_j - \frac{n_{ij}}{r_i\,s_j}
\end{equation*}
has row totals $u$ and column totals $v$. Starting from $s = 1$, it alternates
\begin{equation*}
r_i = \frac{u_i + \sqrt{u_i^2 + 4\,p_i\,n_i}}{2\,p_i}, \qquad
p_i = \sum_j p_{ij}\,s_j, \qquad n_i = \sum_j n_{ij} / s_j
\end{equation*}
($r_i = -n_i / u_i$ when $p_i = 0$), and the same for $s$ with the roles of rows
and columns exchanged, until every row and column total is within the
tolerance of its target. Without negative cells this is RAS. A factor that is
not positive and finite means that the targets cannot be met with the signs of
the cells, and \cmd{cge_balance} stops.

\subsection{Cross entropy}\label{sec:ce}

Only some totals are known, or none: the totals of the other accounts are
estimated with the cells. $X$ minimises $D(X, \bar X)$ subject to
\begin{equation*}
\sum_j x_{kj} - \sum_i x_{ik} = 0 \quad \text{for every account } k,
\qquad
\sum_j x_{kj} = t_k \quad \text{for the accounts whose totals are known.}
\end{equation*}
The solution has the form
\begin{equation*}
x_{ij} = p_{ij}\,e^{a_{ij}} - n_{ij}\,e^{-a_{ij}}, \qquad a_{ij} = \mu_i - \mu_j + \varrho_i,
\end{equation*}
with a multiplier $\mu_k$ for the balance of each account and $\varrho_k$ for
each known total ($\varrho_k = 0$ for the other accounts). The multipliers
minimise the convex dual
\begin{equation*}
\varphi(\mu, \varrho) = \sum_{ij} p_{ij}\,e^{a_{ij}} + \sum_{ij} n_{ij}\,e^{-a_{ij}}
  + \sum_k \mu_k\,c_k + \sum_{k\ \text{known}} \varrho_k\,(f_k - t_k),
\end{equation*}
where $c_k$ is the fixed cells' row total less their column total and $f_k$
their row total. Its gradient is the vector of balances and of the
differences from the known totals, and its Hessian is
$H = \sum_{ij} |x_{ij}|\,d_{ij}\,d_{ij}'$ with $d_{ij} = \partial a_{ij} /
\partial(\mu, \varrho)$. Newton's method with a backtracking line search
minimises $\varphi$ until every balance and known total is within the
tolerance. Adding a constant to all $\mu$ in a connected set of accounts changes
nothing, so one $\mu$ in each set is fixed at zero; an account without cells
that may change must already balance. If the iterations do not converge, no
balanced SAM with the same pattern of signs and zeros may exist.

When every account's total is known, the solution is the RAS solution, with
$r_i = e^{\mu_i + \varrho_i}$ and $s_j = e^{-\mu_j}$; the test suite uses this
identity. This is not the column-coefficient cross entropy of
\citet{robinson2001}, which measures the distance between the column shares of
$X$ and $\bar X$ rather than between the cells; that method is not
implemented.

\subsection{Tolerance and results}

Convergence is measured against the largest absolute cell of $\bar X$: every
balance and target is met to \cmd{tolerance()} times $\max(1, \max_{ij}
|\bar x_{ij}|)$, by default $10^{-10}$, well within the $10^{-8}$ of the balance
check of \cmd{cge_map}. \cmd{cge_balance} reports the largest imbalance before
and after, the number of cells changed, the largest relative change of a cell,
the SAM total before and after, and the cells with the largest changes. The
balanced SAM is an estimate: which totals are known, which cells are fixed and
which method is used should be stated with any results.

\section{Accounting multipliers}\label{sec:multipliers}

\manref{cge_multiplier} computes the accounting multipliers of the mapped SAM
\citep{pyatt1979}: the effects of an injection into one account on the incomes
of all the accounts whose spending the SAM explains, with prices fixed. It uses
the SAM after the layout rewrites of section~\ref{sec:layouts}, and needs no
calibration.

The accounts are divided into endogenous accounts $n$ and exogenous accounts.
By default the endogenous accounts are the activities, commodities and margin
accounts (the production block, $b = 1$), the factors ($b = 2$), and the
households and enterprises (the institutions block, $b = 3$); the government,
savings--investment, the rest of the world, the tax accounts and the health
insurance fund are exogenous. With $t_j$ the column total of account $j$, the
average expenditure propensities among endogenous accounts are
\begin{equation*}
A_{ij} = \SAM(i, j) / t_j, \qquad i, j \in n,
\end{equation*}
(zero for an account without payments), and the incomes of the endogenous
accounts satisfy $y = A\,y + x$, where $x$ is their income from exogenous
accounts. Hence
\begin{equation*}
y = M\,x, \qquad M = (I - A)^{-1}.
\end{equation*}
Column $j$ of $M$ is the change in the incomes of every endogenous account for
a unit injection into $j$. For each injection the command reports the total
multiplier $\sum_i M_{ij}$ and its parts in gross output (the activities),
value added (the factors) and household income (the households). $I - A$ is
invertible when every chain of payments among the endogenous accounts
eventually leaks to an exogenous account (taxes, savings, imports).

\subsection{Decomposition}\label{sec:decomposition}

With the three blocks ordered as above, $A$ has a circular structure: factors
are paid by production ($A_{21}$), institutions by factors ($A_{32}$), and
production by institutions ($A_{13}$), with transfers within production
($A_{11}$), within factors ($A_{22}$) and within institutions ($A_{33}$). With
$\tilde A$ the block-diagonal part of $A$ and $A^* = (I - \tilde A)^{-1}(A -
\tilde A)$, \citet{pyatt1979} write
\begin{equation*}
M = M_3\,M_2\,M_1, \qquad M_1 = (I - \tilde A)^{-1}, \quad M_2 = I + A^* + A^{*2},
\quad M_3 = (I - A^{*3})^{-1},
\end{equation*}
the transfer, open-loop and closed-loop multipliers, and additively
\begin{equation*}
M = I + (M_1 - I) + (M_2 - I)\,M_1 + (M_3 - I)\,M_2\,M_1 .
\end{equation*}
The three terms after $I$ are the transfer effects within the injected
account's block, the open-loop effects on the other blocks, and the closed-loop
effects of the circular flow back to the injected block. The decomposition
requires the blocks outside the circle ($A_{12}$, $A_{23}$, $A_{31}$) to be
zero; \cmd{decompose} stops with an error, listing the largest flows that break
the structure, if they are not.

\subsection{Injections and use}

\cmd{inject()} gives a vector $x$ of injections into endogenous accounts, for
example more government purchases of a commodity, and the command reports
$M\,x$: the change in gross output by activity, in factor incomes and in
household incomes, and their totals.

Accounting multipliers assume fixed prices, unused capacity in every sector and
average propensities equal to marginal ones; they are upper bounds on the
effects that a CGE model, with its supply constraints and price changes, would
find. They complement the CGE model; they do not replace it.

\section{Structural path analysis}\label{sec:spa}

\manref{cge_spa} decomposes one accounting multiplier into the paths along
which an injection travels through the SAM \citep{defourny1984}. The accounts
are the poles of a graph, with an arc from pole $i$ to pole $j$ when $j$
receives part of $i$'s spending, with intensity $A_{ji}$. An elementary path
$p$ from an origin $i$ to a destination $j$ is a sequence of arcs that passes
through no pole twice.
\begin{itemize}
\item The \textbf{direct influence} of $i$ on $j$ along $p$ is the product of
the intensities of its arcs, $I^D_p = A_{k_1 i}\,A_{k_2 k_1} \cdots A_{j k_m}$:
the effect along the path with every other account held fixed.
\item The \textbf{path multiplier} $M_p$ adds the feedback of the loops that
start and end at the poles of $p$. With $P$ the set of poles of $p$,
\begin{equation*}
M_p = \frac{\det\bigl((I - A)_{-P,-P}\bigr)}{\det(I - A)} = \det\bigl(M_{P,P}\bigr),
\end{equation*}
where $(I - A)_{-P,-P}$ leaves out the rows and columns of $P$ and $M_{P,P}$ is
the block of $M$ on $P$; the two are equal by Jacobi's identity for
complementary minors, and the second needs only a small determinant.
\item The \textbf{total influence} along $p$ is $I^T_p = I^D_p\,M_p$.
\item The \textbf{global influence} of $i$ on $j$ is the multiplier $M_{ji}$,
and it is the sum of the total influences along all elementary paths from $i$
to $j$ ($i \ne j$).
\end{itemize}
The number of elementary paths grows quickly with the size of the SAM, so the
search stops at paths of \cmd{maxlength()} arcs (default 6) and drops a path,
and every extension of it, once its direct influence falls below
\cmd{threshold()} times the global influence in absolute value (default 0.001).
With intensities no larger than one, as when a SAM has no negative cells,
extensions can only lower the direct influence, so no path above the threshold
is missed. The share of the global influence explained by the paths found is
reported.
